% !TEX root = main.tex
\chapter{Αλγόριθμοι Δειγματοληψίας Langevin}
\label{chap:algorithms}

Το κεφάλαιο αυτό περιγράφει τους τρεις αλγορίθμους που συγκρίνονται στη
εργασία αυτή ως διακριτοποιήσεις των SDE Langevin της
Ενότητας~\ref{sec:bg-langevin}. Ξεκινάμε με τις δύο overdamped παραλλαγές,
ULA και MALA (Ενότητες~\ref{sec:alg-ula}--\ref{sec:alg-mala}), και στη
συνέχεια περνάμε στο underdamped (kinetic) σχήμα BAOAB
(Ενότητα~\ref{sec:alg-kinetic}). Η Ενότητα~\ref{sec:alg-momentum-role}
εξηγεί τον ρόλο της ορμής και τους συμβιβασμούς μεταξύ των τριών αλγορίθμων.

\section{Overdamped Langevin: ULA}
\label{sec:alg-ula}

Η overdamped SDE Langevin \eqref{eq:overdamped-bg} διακριτοποιείται με βήμα
Euler--Maruyama μεγέθους $\eta > 0$, από την οποία προκύπτει η αναδρομή
\begin{equation}
  X_{k+1} \;=\; X_k - \eta\, \nabla U(X_k) + \sqrt{2\eta}\, \xi_k,
  \qquad \xi_k \overset{iid}{\sim} \mathcal{N}(0, I_d),
  \label{eq:ula}
\end{equation}
που είναι ο \emph{Unadjusted Langevin Algorithm} (ULA). Η εξίσωση \eqref{eq:ula}
είναι ο απλούστερος δυνατός δειγματολήπτης Langevin: μία αξιολόγηση κλίσης ανά
επανάληψη, μία Γκαουσιανή διέγερση, χωρίς έλεγχο αποδοχής/απόρριψης.

Το τίμημα αυτής της απλότητας είναι μια μεροληψία. Η αλυσίδα Markov που ορίζεται
από το \eqref{eq:ula} έχει αναλλοίωτη κατανομή $\pi_\eta$ που διαφέρει από την
κατανομή-στόχο $\pi$ κατά σφάλμα $\mathcal{O}(\eta)$. Η διακριτοποίηση έχει
αδύναμη τάξη ένα ως προς $\eta$ \citep{kloeden1992numerical}.
Ο \citet{roberts1996exponential} απέδειξε εκθετική εργοδικότητα του
\eqref{eq:ula} υπό παραδοχές έντονης κυρτότητας για το $U$, και ο
\citet{durmus2017nonasymptotic} παρέσχε ποσοτικά μη-ασυμπτωτικά φράγματα στην
απόσταση ολικής μεταβολής μεταξύ $\pi_\eta$ και $\pi$ ως συνάρτηση του $\eta$,
του μήκους αλυσίδας $n$ και ιδιοτήτων του $U$. Για τους σκοπούς μας, το κρίσιμο
χαρακτηριστικό είναι ότι η στατική κατανομή του ULA έχει μεροληψία που δεν
εξαφανίζεται καθώς το μήκος αλυσίδας $n$ αυξάνεται: μεγαλύτερες αλυσίδες
συγκλίνουν στη λανθασμένη κατανομή.

Αλγοριθμικά, ο ULA παρατίθεται στον Αλγόριθμο~\ref{alg:ula}. Η υλοποίηση που
χρησιμοποιείται στην εργασία βρίσκεται στο
\texttt{samplers/overdamped/ula/ula\_runner.py}.

\begin{algorithm}[ht]
\caption{Unadjusted Langevin Algorithm (ULA)}
\label{alg:ula}
\begin{algorithmic}[1]
\Require βήμα $\eta > 0$, αριθμός βημάτων $N$, αρχική κατάσταση $x_0 \in \mathbb{R}^d$, κλίση $\nabla U$
\State $x \gets x_0$
\For{$k = 1, 2, \ldots, N$}
  \State Δειγματοληψία $\xi_k \sim \mathcal{N}(0, I_d)$
  \State $x \gets x - \eta\, \nabla U(x) + \sqrt{2\eta}\, \xi_k$
  \State Καταγραφή $X_k \gets x$
\EndFor
\State \Return $X_1, \ldots, X_N$
\end{algorithmic}
\end{algorithm}

Μια πρακτική επιπλοκή του ULA είναι η ευαισθησία του στο $\eta$: η επιλογή πολύ
μικρού $\eta$ αυξάνει τον ολοκληρωμένο χρόνο αυτοσυσχέτισης, ενώ η επιλογή
πολύ μεγάλου $\eta$ αυξάνει τη μόνιμη μεροληψία. Το βέλτιστο $\eta$ εξαρτάται από
την τοπική καμπυλότητα του $U$, η οποία σε κατανομές-στόχους με βαριές ουρές
ποικίλλει απότομα στον χώρο καταστάσεων. Έτσι δεν υπάρχει βήμα που να αποδίδει
ικανοποιητικά παντού. Βλ.\ Ενότητες~\ref{sec:grid-study} και \ref{sec:tail-correlation} για τις
εμπειρικές εκδηλώσεις αυτού του φαινομένου.

\section{Overdamped Langevin με Διόρθωση Metropolis: MALA}
\label{sec:alg-mala}

Ο \emph{Metropolis-Adjusted Langevin Algorithm} (MALA) αποκαθιστά την ακρίβεια
σε οποιοδήποτε βήμα συνδυάζοντας την πρόταση ULA με βήμα αποδοχής/απόρριψης
Metropolis--Hastings. Ξεκινώντας από $X_k = x$, ο MALA προτείνει
\begin{equation}
  X' \;=\; x - \eta\, \nabla U(x) + \sqrt{2\eta}\, \xi_k,
  \qquad \xi_k \sim \mathcal{N}(0, I_d),
\end{equation}
και αποδέχεται με πιθανότητα $\alpha(x, X')$ όπου
\begin{equation}
  \alpha(x, x') \;=\; \min\!\left(1, \;
    \frac{\pi(x')\, q(x \mid x')}{\pi(x)\, q(x' \mid x)}
  \right),
  \label{eq:mh-prob}
\end{equation}
και η κατανομή πρότασης $q$ είναι Γκαουσιανή με μέσο $x - \eta\nabla U(x)$
και συνδιακύμανση $2\eta I_d$:
\begin{equation}
  \log q(x' \mid x) \;=\; -\frac{1}{4\eta} \left\| x' - x + \eta \nabla U(x) \right\|^2
    \;+\; \text{const}
  \label{eq:mala-proposal}
\end{equation}
Η διόρθωση Metropolis επιβάλλει λεπτομερή ισορροπία ως προς την $\pi$, οπότε
η αλυσίδα είναι ακριβής για οποιοδήποτε $\eta > 0$
\citep{roberts1998optimal, robert2004monte}.

Το κόστος της διόρθωσης είναι μη-τετριμμένο. Κάθε επανάληψη απαιτεί (i) μια
επιπλέον αξιολόγηση του $\pi$ (ή $U$) στην πρόταση, και (ii) η αποδοχή δεν
είναι εγγυημένη: η αλυσίδα σπαταλά υπολογιστικούς πόρους σε απορριπτόμενες
προτάσεις. Η θεωρία βέλτιστης κλίμακας του \citet{roberts1998optimal} αποδεικνύει
ότι το ποσοστό αποδοχής που μεγιστοποιεί την αναμενόμενη τετραγωνική μετατόπιση
ανά μονάδα υπολογισμού είναι περίπου $0.574$ στο υψηλής-διάστασης όριο για
κατανομές-στόχους Γκαουσιανής μορφής γινομένου. Αυτό αποτελεί τη βάση του
φίλτρου $[0.2, 0.95]$ που εφαρμόζεται στις διαμορφώσεις MALA κατά την επιλογή
των βέλτιστων παραμέτρων στην Ενότητα~\ref{sec:grid-study}.

Ο MALA συνοψίζεται στον Αλγόριθμο~\ref{alg:mala}.

\begin{algorithm}[ht]
\caption{Metropolis-Adjusted Langevin Algorithm (MALA)}
\label{alg:mala}
\begin{algorithmic}[1]
\Require βήμα $\eta > 0$, αριθμός βημάτων $N$, αρχική κατάσταση $x_0$, δυναμικό $U$ και κλίση $\nabla U$
\State $x \gets x_0$
\For{$k = 1, 2, \ldots, N$}
  \State Δειγματοληψία $\xi_k \sim \mathcal{N}(0, I_d)$
  \State $x' \gets x - \eta\, \nabla U(x) + \sqrt{2\eta}\, \xi_k$ \Comment{Πρόταση}
  \State Υπολογισμός $\log \alpha = U(x) - U(x') + \log q(x \mid x') - \log q(x' \mid x)$
  \State Δειγματοληψία $u \sim \mathrm{Uniform}(0,1)$
  \If{$\log u < \log \alpha$}
    \State $x \gets x'$ \Comment{Αποδοχή}
  \EndIf
  \State Καταγραφή $X_k \gets x$
\EndFor
\State \Return $X_1, \ldots, X_N$
\end{algorithmic}
\end{algorithm}

Για κατανομές με βαριές ουρές, το ποσοστό αποδοχής του MALA έχει μια
ενδιαφέρουσα συμπεριφορά στην οποία επανερχόμαστε στο
Κεφάλαιο~\ref{chap:experiments}. Σε σταθερό βήμα $\eta$, ο MALA αποδέχεται
\emph{συχνότερα} στη Student-$t$ παρά στη Γκαουσιανή, επειδή η λογαριθμική
πυκνότητα των βαρύτερων ουρών είναι πιο επίπεδη και ο λόγος Metropolis λιγότερο
κεντρωμένος. Στην Cauchy υψηλής διάστασης, το ποσοστό αποδοχής πέφτει
σημαντικά σε σύγκριση με τη χαμηλής διάστασης περίπτωση --- η πιθανότητα αποδοχής
οποιουδήποτε μεγάλου άλματος κλιμακώνεται ως $\exp(-(d+1)/2 \cdot \log(1 +
\|x'\|^2))$, δηλαδή ως \emph{δύναμη} του $\|x'\|$ και όχι ως εκθετική, οπότε ο
λόγος Metropolis είναι πολύ πιο ευαίσθητος στη γεωμετρία της πρότασης απ' ό,τι
για κατανομές-στόχους με Γκαουσιανές ουρές. Το αποτέλεσμα είναι ότι το
παραγωγικό βήμα στην Cauchy σε μέτρια διάσταση είναι αναγκαστικά μεγαλύτερο από
ό,τι στη Γκαουσιανή, αλλά το αντίστοιχο ποσοστό αποδοχής είναι χαμηλότερο.

\section{Underdamped Langevin και Splitting Ολοκληρωτές}
\label{sec:alg-kinetic}

Η underdamped SDE Langevin \eqref{eq:underdamped-bg} εισάγει μια μεταβλητή ορμής
$V_t \in \mathbb{R}^d$ παράλληλα με τη θέση $X_t$ και προσθέτει έναν όρο τριβής
$-\gamma V_t \, dt$ στην ταχύτητα. Η $x$-περιθώρια κατανομή του κοινού αναλλοίωτου
μέτρου είναι η κατανομή-στόχος $\pi$, οπότε η δειγματοληψία από την $\pi$
ανάγεται στην ολοκλήρωση της κοινής SDE και απόρριψη της συνιστώσας ορμής. Το
kinetic καθεστώς αντιστοιχεί σε μικρό $\gamma$ (χαμηλή τριβή, βαλλιστική κίνηση
μέσα στο support). Το overdamped όριο ανακτάται καθώς $\gamma \to \infty$ υπό
κατάλληλη επανακλίμακα χρόνου.

\subsection{Γιατί Splitting;}
\label{sec:alg-splitting-motivation}

Μια αφελής διακριτοποίηση Euler--Maruyama του \eqref{eq:underdamped-bg} υποφέρει
από τα ίδια προβλήματα ακρίβειας πρώτης τάξης με τον ULA. Μια πιο εξελιγμένη
προσέγγιση είναι να \emph{χωρίσουμε} τη ΣΔΕ σε τμήματα που μπορεί το καθένα να
ολοκληρωθεί ακριβώς ή με υψηλότερη ακρίβεια, και να συνθέσουμε τις λύσεις
\citep{leimkuhler2015molecular}. Γράφοντας τον underdamped γεννήτορα ως
$\mathcal{L} = A + B + O$ με
\begin{align}
  A:\quad & dX_t = V_t\, dt, \qquad dV_t = 0,                              \tag{ολίσθηση στο $X$} \\
  B:\quad & dX_t = 0,         \qquad dV_t = -\nabla U(X_t)\, dt,           \tag{δύναμη στο $V$} \\
  O:\quad & dX_t = 0,         \qquad dV_t = -\gamma V_t\, dt + \sqrt{2\gamma}\, dW_t, \tag{Ornstein--Uhlenbeck στο $V$}
\end{align}
κάθε τμήμα μπορεί να επιλυθεί ακριβώς. Ειδικότερα, το βήμα $O$ είναι μια διαδικασία
Ornstein--Uhlenbeck με αναλυτική ενημέρωση
\begin{equation}
  V_{t+\eta} \;=\; e^{-\gamma \eta}\, V_t \;+\; \sqrt{1 - e^{-2\gamma\eta}}\, \xi,
  \qquad \xi \sim \mathcal{N}(0, I_d),
  \label{eq:ou-step}
\end{equation}
που διατηρεί ακριβώς την περιθώρια $\mathcal{N}(0, I_d)$ της ορμής για οποιοδήποτε
βήμα $\eta$ \citep{leimkuhler2013bao}. Τα υπόλοιπα τμήματα $A$ και $B$ είναι
ντετερμινιστικά και μπορούν να ολοκληρωθούν ακριβώς με ενημερώσεις μισού βήματος.

\subsection{Ο Ολοκληρωτής BAOAB}
\label{sec:alg-baoab}

Το σχήμα \emph{BAOAB} του \citet{leimkuhler2013bao} συνθέτει τις στοιχειώδεις
ροές με σειρά $B \to A \to O \to A \to B$, κάθε μία με το μισό ονομαστικό βήμα
για τα $A$ και $B$ και ολόκληρο το βήμα για το $O$.  Ο ολοκληρωτής αυτός είναι ιδιαίτερα ακριβής για Γκαουσιανές κατανομές-στόχους καθώς το σφάλμα κλιμακώνεται ως $\mathcal{O}(\eta^3)$ αντί για $\mathcal{O}(\eta)$ του απλού Euler και αποτελεί το σχήμα αναφοράς για kinetic δειγματολήπτες Langevin στη μοριακή δυναμική \citep{leimkuhler2015molecular}.


Ένα βήμα BAOAB ξεκινώντας από $(x, v)$ εκτελείται ως εξής:
\begin{align}
  \text{B}:\quad & v \leftarrow v - \tfrac{\eta}{2}\, \nabla U(x), \\
  \text{A}:\quad & x \leftarrow x + \tfrac{\eta}{2}\, v, \\
  \text{O}:\quad & v \leftarrow e^{-\gamma \eta}\, v + \sqrt{1 - e^{-2\gamma\eta}}\, \xi, \qquad \xi \sim \mathcal{N}(0, I_d), \\
  \text{A}:\quad & x \leftarrow x + \tfrac{\eta}{2}\, v, \\
  \text{B}:\quad & v \leftarrow v - \tfrac{\eta}{2}\, \nabla U(x).
\end{align}
Αυτή η ακολουθία ενημερώσεων παράγει τη νέα κατάσταση $(X_{k+1}, V_{k+1})$.
Όπως ο ULA, το σχήμα BAOAB δεν περιλαμβάνει διόρθωση Metropolis, οπότε η στατική
κατανομή του διαφέρει από την κατανομή-στόχο κατά σφάλμα διακριτοποίησης. Η
δομή του σφάλματος είναι ωστόσο καλύτερα ελεγχόμενη από εκείνη του ULA.
Συγκεκριμένα, σε γενικό λείο δυναμικό η περιθώρια κατανομή θέσης διαφέρει
από την κατανομή-στόχο κατά $\mathcal{O}(\eta^2)$ σε ολική μεταβολή
\citep{leimkuhler2013bao, leimkuhler2013rational}, σε αντίθεση με το
$\mathcal{O}(\eta)$ του ULA. Σε τετραγωνικό δυναμικό η περιθώρια κατανομή
θέσης ανακτάται \emph{ακριβώς} (οι ροπές ταιριάζουν με τις τιμές-στόχους έως
στρογγυλοποίηση), κάτι που ο \citet{leimkuhler2013bao} αποκαλεί
\emph{ρυθμιστική υπερ-σύγκλιση} του BAOAB. Αυτό σημαίνει ότι στην πράξη,
σε σχεδόν-τετραγωνικά δυναμικά, ο BAOAB έχει πολύ μικρότερη μεροληψία
από τον ULA για το ίδιο βήμα $\eta$ \citep{leimkuhler2013rational}.

Ο ψευδοκώδικας του δειγματολήπτη BAOAB εμφανίζεται ως Αλγόριθμος~\ref{alg:baoab}.
Η υλοποίηση που χρησιμοποιείται στην εργασία βρίσκεται στο
\texttt{samplers/underdamped/kinetic/splitting/baoab/kinetic\_baoab\_runner.py}.

\begin{algorithm}[ht]
\caption{Ολοκληρωτής BAOAB για underdamped Langevin}
\label{alg:baoab}
\begin{algorithmic}[1]
\Require βήμα $\eta>0$, τριβή $\gamma > 0$, αριθμός βημάτων $N$, αρχική κατάσταση $(x_0, v_0)$, κλίση $\nabla U$
\State $x, v \gets x_0, v_0$
\State Προϋπολογισμός $a \gets \exp(-\gamma \eta)$ και $s \gets \sqrt{1 - a^2}$
\For{$k = 1, 2, \ldots, N$}
  \State Δειγματοληψία $\xi_k \sim \mathcal{N}(0, I_d)$
  \State $v \gets v - \tfrac{\eta}{2}\, \nabla U(x)$ \Comment{B}
  \State $x \gets x + \tfrac{\eta}{2}\, v$ \Comment{A}
  \State $v \gets a\, v + s\, \xi_k$ \Comment{O}
  \State $x \gets x + \tfrac{\eta}{2}\, v$ \Comment{A}
  \State $v \gets v - \tfrac{\eta}{2}\, \nabla U(x)$ \Comment{B}
  \State Καταγραφή $X_k \gets x$
\EndFor
\State \Return $X_1, \ldots, X_N$
\end{algorithmic}
\end{algorithm}

Υπολογιστικά, κάθε βήμα BAOAB απαιτεί \emph{δύο} αξιολογήσεις κλίσης (στις δύο
ενημερώσεις B), έναντι μιας ανά βήμα για τον ULA και μιας ανά αποδεκτή πρόταση
για τον MALA. Το κόστος κλίσης είναι επομένως περίπου $1.5\times$--$2\times$
εκείνο του ULA, χωρίς να υπολογίζονται οι ντετερμινιστικές ενημερώσεις A που
είναι φτηνές.

\section{Ο Ρόλος της Ορμής}
\label{sec:alg-momentum-role}

Το εμπειρικό πλεονέκτημα της underdamped διατύπωσης έναντι της overdamped ---
εμφανές στα αποτελέσματα του Κεφαλαίου~\ref{chap:experiments} --- μπορεί να
εξηγηθεί από τον τρόπο που η ορμή μεταφέρει πληροφορία από βήμα σε βήμα.

Στο overdamped σχήμα, η θέση ενημερώνεται στο βήμα $k$ αποκλειστικά βάσει
της κλίσης $\nabla U(X_k)$ και μιας νέας Γκαουσιανής διέγερσης $\xi_k$. Δεν
υπάρχει μνήμη της πρόσφατης ιστορίας της τροχιάς. Τα διαδοχικά βήματα
θέσης είναι επομένως σχεδόν ανεξάρτητα μεταξύ τους (πέρα από τον όρο
κλίσης), οπότε η αλυσίδα μετακινείται διαχυτικά, σε χρονική κλίμακα
$\sqrt{n\eta}$ για τροχιά $n$ βημάτων.

Στο underdamped σχήμα, η μεταβλητή ορμής παίζει τον ρόλο μιας κατάστασης που
συσσωρεύει Γκαουσιανές διεγέρσεις και φθίνει γεωμετρικά μέσω του όρου τριβής
$-\gamma V \, dt$. Με τριβή $\gamma \eta \ll 1$ (kinetic καθεστώς), η ορμή
διατηρείται κατά προσέγγιση σε κάθε βήμα O, και τα διαδοχικά βήματα
θέσης είναι \emph{θετικά συσχετισμένα}, οπότε η αλυσίδα κινείται βαλλιστικά και
όχι διαχυτικά, σε χρονική κλίμακα $n\eta$. Αυτός είναι ο ίδιος μηχανισμός που
διέπει το Hamiltonian Monte Carlo \citep{neal2011mcmc}, του οποίου το underdamped
Langevin μπορεί να θεωρηθεί στοχαστικό ανάλογο.

Οι εμπειρικές συνέπειες είναι πιο εμφανείς σε κατανομές-στόχους με βαριές ουρές:
για να φτάσει η αλυσίδα στις βαθιές ουρές της Cauchy γύρω στο $\|x\| \sim 30$,
χρειάζεται μεγάλη ολική μετατόπιση που το overdamped σχήμα μπορεί να επιτύχει
μόνο διαχυτικά, αλλά το underdamped σχήμα μπορεί να παρέχει σε μια ενιαία
διέλευση κινούμενη υπό ορμή. Τα trace plots στην Ενότητα~\ref{sec:case-studies} δείχνουν
αυτή τη διαφορά καθαρά για ULA και BAOAB στην ίδια κατανομή-στόχο Cauchy.

Ο συμβιβασμός βρίσκεται στην παράμετρο τριβής $\gamma$. Μικρό $\gamma$ μεγιστοποιεί
το πλεονέκτημα βαλλιστικής μεταφοράς αλλά μειώνει τις $\pi$-διατηρητικές
ιδιότητες του βήματος O (η αλυσίδα εξερευνά υπερβολικά ελεύθερα, και ο ρυθμός
ανάμειξης προς το αναλλοίωτο μέτρο μπορεί να μειωθεί). Μεγάλο $\gamma$ ανακτά
το overdamped όριο και χάνει το kinetic πλεονέκτημα. Η εμπειρική διερεύνηση
στην Ενότητα~\ref{sec:grid-study} επιβεβαιώνει ότι στις κατανομές-στόχους μας
το κατώτερο άκρο του εξερευνούμενου πλέγματος $\gamma$ προτιμάται σταθερά ---
κάτι που επιβεβαιώνει ότι το kinetic πλεονέκτημα ισχύει και για βαριές ουρές.

\section{Σύγκριση Αλγορίθμων με Μια Ματιά}
\label{sec:alg-summary}

Ο Πίνακας~\ref{tab:alg-summary} συνοψίζει τις αλγοριθμικές και υπολογιστικές
ιδιότητες των τριών δειγματοληπτών.

\begin{table}[ht]
\centering
\small
\renewcommand{\arraystretch}{1.15}
\begin{tabular}{@{}p{0.27\textwidth}>{\raggedright\arraybackslash}p{0.18\textwidth}>{\raggedright\arraybackslash}p{0.20\textwidth}>{\raggedright\arraybackslash}p{0.27\textwidth}@{}}
\toprule
                                  & ULA      & MALA              & BAOAB                  \\
\midrule
SDE                               & overdamped & overdamped       & underdamped            \\
Βοηθητική μεταβλητή              & καμία    & καμία             & ορμή $V$               \\
Διόρθωση Metropolis               & όχι      & ναι               & όχι                    \\
Ακριβής για $\eta > 0$;           & όχι      & ναι               & όχι                    \\
Μεροληψία        & $\mathcal{O}(\eta)$ & $0$    & $\mathcal{O}(\eta^2)$ (γενική περίπτωση). Ακριβής σε τετραγωνικό δυναμικό \\
Αξιολογήσεις κλίσης / βήμα      & $1$      & $1$ (+ απορριπτ.) & $2$                    \\
Αποδοχή/Απόρριψη                  & όχι      & ναι               & όχι                    \\
Ρυθμιζόμενες παράμετροι          & $\eta$   & $\eta$            & $\eta, \gamma$         \\
Ενδείκνυται για\ldots            & φτηνές διερευνητικές εκτελέσεις & μέτριες ουρές, μικρό $d$ & ελαφριές/μέτριες ουρές οποιοδήποτε $d$. Βαριές ουρές \\
\bottomrule
\end{tabular}
\caption{Αλγοριθμική σύγκριση των τριών δειγματοληπτών Langevin που μελετώνται
στην εργασία. Η γραμμή «Ενδείκνυται για» προαναγγέλλει τα εμπειρικά
ευρήματα του Κεφαλαίου~\ref{chap:experiments}.}
\label{tab:alg-summary}
\end{table}

Στο Κεφάλαιο~\ref{chap:experiments} βλέπουμε αυτές τις διαφορές με αριθμούς
και συγκεκριμένα παραδείγματα.
